\subsection{Gradiente descendiente}
\label{sec:gradiente_descendiente}

El método de gradiente descendiente consiste en actualizar los parámetros en la dirección contraria al gradiente de la función de costo ($\nabla_{\theta} L(y, \hat{y})$), es decir, en la dirección que reduce el error \cite{Cauchy1847}. La actualización de los parámetros se realiza como se muestra en la ecuación \ref{eq:gradient_descent}, donde $\eta$ es la tasa de aprendizaje, un hiperparámetro, que determina el tamaño del paso en cada iteración.

\begin{equation}
    \label{eq:gradient_descent}
    \theta_{n+1} = \theta_n - \eta \nabla_{\theta} L(y, \hat{y})
\end{equation}

El gradiente de una función define la dirección de mayor incremento de la función en un punto determinado. Por lo tanto, al restar el gradiente en la actualización de los parámetros, se avanza en la dirección de mayor decremento de la función de costo, buscando minimizarla iterativamente.

Es posible también caer en mínimos locales de la función de costo, dependiendo de la forma de la función y del punto de inicio. Sin embargo, en muchos casos prácticos, el gradiente descendiente puede encontrar soluciones suficientemente buenas, incluso si no son óptimas globales \cite{Bishop2006}.

En esta búsqueda iterativa, la elección de la tasa de aprendizaje $\eta$ es crucial. Una tasa de aprendizaje demasiado alta puede provocar que el algoritmo diverja, mientras que una tasa demasiado baja puede resultar en una convergencia lenta \cite{Bishop2006}.
